\documentclass[leqno]{jsarticle}
\usepackage{amssymb}
\usepackage{amsthm}
\usepackage{amsmath}
\usepackage{comment}
\usepackage[all]{xy}
%定理環境 thmはあまり使わずpropをなるべく使う。
%主張は基本的に証明の中で使い、そのため番号を付けない。
%注意環境を付けてみた。
\newtheorem{thm}{定理}
\newtheorem{prop}[thm]{命題}
\newtheorem{cor}[thm]{系}
\newtheorem{lem}[thm]{補題}
\newtheorem{df}[thm]{定義}
\newtheorem{exam}[thm]{例}
\newtheorem{fact}[thm]{事実}
\newtheorem*{claim}{主張}
\newtheorem*{cau}{注意}
\renewcommand{\proofname}{{\bf 証明}}
%矢印たち。
%\toは\rightarrowと同じ。\getsは\leftarrowと同じ。同値は\iff
%上向き矢はuparrow逆はdownarrow
\newcommand{\To}{\Longrightarrow}
\newcommand{\Gets}{\LongLeftarrow}
%黒板太字
\newcommand{\nn}{\mathbb{N}}
\newcommand{\zz}{\mathbb{Z}}
\newcommand{\qq}{\mathbb{Q}}
\newcommand{\rr}{\mathbb{R}}
\newcommand{\cc}{\mathbb{C}}
%無限大
\newcommand{\mgn}{\infty}
%差集合は\setminusで統一する。等号の否定は\neq
%真部分集合は\subsetneqq　偏微分は\partial
%ディスプレイスタイル。\dfracでディスプレイ分数
\newcommand{\dpst}{\displaystyle}
\newcommand{\ep}{\varepsilon}
\newcommand{\qt}[2]{#1/\! #2}
\newcommand{\al}{\alpha}
\newcommand{\be}{\beta}
\newcommand{\la}{\lambda}
\newcommand{\La}{\Lambda}
\newcommand{\ga}{\gamma}
\newcommand{\hk}{\hookrightarrow}
\newcommand{\pt}{\partial}
%空集合は\Oを使う！！！！！！！！！！！！

\newcommand{\main}{X\underset{f}{\cup} Y}
\title{接着、帰納極限、CW複体}
\author{電波通信}
\date{\today}
			\begin{document}
\maketitle
\section{接着空間}
\begin{df}
全射連続写像$f:X\to Y$が等化写像であるとは
\[
\mbox{$U$が$Y$の開集合}\iff\mbox{$f^{-1}(U)$が$X$の開集合}
\]
が成り立つときに言う。このとき同時に
\[
\mbox{$F$が$Y$の閉集合}\iff\mbox{$f^{-1}(F)$が$X$の閉集合}
\]
が成立する。
\end{df}
例えば位相空間$X$上の同値関係$R$から誘導される商写像$p:X\to \qt{X}{R}$は等化写像となっている。またこのことのある意味での逆が成り立つ。
\begin{prop}\label{syo}
等化写像$f:X\to Y$について$X$上の同値関係$R$を
\[
xRy\overset{def}{\iff}f(x)=f(y)
\]
と定義すると
\[
\qt{X}{R}\approx Y
\]
である。
\end{prop}
\begin{proof}$R$の同値類は$f^{-1}(y)$という形をしている。$xRy\To f(x)=f(y)$から下の可換図式を充たす連続関数$F:\qt{X}{R}\to Y$が存在する。
\[\xymatrix{
& X\ar@{->}[r]^{\pi} \ar@{->}[rd]_f& \qt{X}{R}\ar@{-->}[d]^F&\\ 
&  & Y&
}\]
$R$の定義から$F$は全単射である。$F$が開集合であることを示そう。$U$を$\qt{X}{R}$の開集合として$O=\pi^{-1}(U)$と置くと$O$は$X$の開集合であり$\pi(O)=U$を充たす。そして
\[
O=\bigcup_{z\in U}f^{-1}(F(z))
\]である。
さて
\[
F(U)=F(\pi(O))=F\circ \pi(O)=f(O)
\]
である。そして
\[
f^{-1}(f(O))=O
\]
が成り立つので$f(O)$は$Y$の開集合で$F$は開写像となる。よって$F$は同相写像
\end{proof}
\begin{df}
$X,Y$を位相空間とし$A\subset X$を$X$の閉集合、$f:A\to Y$を連続写像とするとき$X$と$Y$を$f$で接着空間
\[
X\underset{f}{\cup} Y
\]
を定義するが、$X$と$Y$の位相的直和空間を$Z=X\sqcup Y$として
\[
\{(a,f(a))\in Z\times Z\ |\ a\in A\}
\]
から生成される最小の$Z$上の最小の同値関係を$R$として
\[
X\underset{f}{\cup} Y=\qt{Z}{R}
\]
と定義する。$X\underset{f}{\cup} Y$の位相は商空間$\qt{Z}{R}$としての商位相を導入する。このとき次の可換図式がある。
\[\xymatrix{
& A \ar@{^{(}->}[r]^{\iota}\ar@{->}[d]_f& X \ar@{->}[d]^{i}&\\ 
& Y \ar@{->}[r]_{j}& X\underset{f}{\cup} Y &
}\]
ここで$i$は写像の合成$X\to X\sqcup Y\to X\underset{f}{\cup} Y$であり、$j$は写像の合成$Y\to X\sqcup Y\to X\underset{f}{\cup} Y$である。
\end{df}
\begin{cau}
$A=\O$の場合もありえる。その場合、$\main=X\sqcup Y$	となる。
\end{cau}
\begin{fact}[この文章では使わない。]
位相空間$Q$と$a:X\to Q,\ b:Y\to Q$が次の可換図式
\[\xymatrix{
& A \ar@{^{(}->}[r]^{\iota}\ar@{->}[d]_f& X \ar@{->}[d]^{a}&\\ 
& Y \ar@{->}[r]_{b}& Q &
}\]
を充たすとき、写像$e:X\underset{f}{\cup} Y\to Q$が存在して次の可換図式を充たす
\[\xymatrix{
& A \ar@{^{(}->}[r]^{\iota}\ar@{->}[d]_f& X \ar@{->}[d]_{i}\ar@{->}[rdd]^a&\\ 
& Y \ar@{->}[r]^{j}\ar@{->}[rrd]_b& X\underset{f}{\cup} Y \ar@{-->}[rd]^e &\\ 
&   &   &  Q&
}\]
要するに$X\underset{f}{\cup} Y$は図式$Y\gets A\hookrightarrow X$のプッシュアウトなのである。この命題は使わないので証明はしないが、商空間の性質を用いれば簡単に示せる。使わないとは言え大事な性質なので言及するだけしておいた。
\end{fact}
\begin{prop}\label{Y}
$Y\to \main$は閉写像でかつ埋め込みである。つまり$Y$は$\main$の閉部分空間とみなせる。
\end{prop}
\begin{proof}
簡単なので省略。
\end{proof}

分離公理との関係：
以下商写像$Z=X\sqcup Y\to X\underset{f}{\cup} Y$を$\pi$と書く。また$z\in Z$の同値類を$[z]$と書くことにする。
\begin{prop}
$X,Y$が$T_1$であるとき、$X\underset{f}{\cup} Y$も$T_1$である。
\end{prop}
\begin{proof}
$\qt{Z}{R}=X\underset{f}{\cup} Y$であるが、$R$による$Z$の同値類は以下のいづれかの形をしている。つまり$X\underset{f}{\cup} Y$の１点$a$の$\pi$による引き戻し$\pi^{-1}(a)$は次のいづれかの集合になる。\footnote{どういう場合に$\{z\}$の形の集合になるか、どういう場合に$f^{-1}(y)\sqcup \{y\}$の形の集合になるか考えてみよ。}
\begin{align*}
&\{z\}\\
 &f^{-1}(y)\sqcup \{y\}
\end{align*}
よってこれらは$X,Y$の$T_1$性と$f$の連続性からいづれも$Z$の閉集合であり、商位相の定義から$X\underset{f}{\cup} Y$の１点は閉集合になる。
\end{proof}

\begin{prop}\label{seiki}
$X,Y$が正規であるとき、$X\underset{f}{\cup} Y$も正規である。
\end{prop}
\begin{proof}
ティーチェの拡張定理が成り立つこと、すなわち任意の閉集合上の$[0,1]$への連続関数が空間に全体に拡張できることと正規性が同値であった。よって閉集合上の関数が拡張できることによって正規性を示す。\par
$F$を$\main$の閉集合とし、$g$を連続関数$g:F\to [0,1]$とする。すると$\pi^{-1}(F)$は$X\sqcup Y$の閉集合であり、
\[
K=\pi^{-1}(F)\cap X,\ L=\pi^{-1}(F)\cap Y
\]
と置けば$K,L$はそれぞれ$X,Y$の閉集合である。$Y$が正規であることから連続関数$g\circ j:L\to[0,1]$は$Y$全体に拡張出来る。これを$\be$と置く。そして$X$上の閉集合$A\cup K$について$h:A\cup K\to[0,1]$を
\[
h(x)=
\begin{cases}
\be \circ f(x) & \text{$(x\in A)$}\\ 
g\circ i(x) & \text{$(x\in K)$}
\end{cases}
\]
と置く。之がちゃんと定義できていることを見よう。$x\in A\cap K$を取ると$xRf(x)$より$f(x)\in \pi^{-1}(F)$で$f(x)\in Y$であるから$f(x)\in L$である。よって$\be$は$g\circ j$の拡張なので$\main$の図式の可換性から
\[
\be \circ f(x)=\be(f(x))=g\circ j(f(x))=g\circ j\circ f(x)=g\circ i\circ \iota(x)=g\circ i(x)
\]
なので$h$はちゃんと$A\cup K$上で定義されており、$A,K$それぞれの上で連続なので$h$は$A\cup K$上で連続である。そこで$X$の正規性から$h$を$X$上までに連続的に拡張し、その関数を$\al$とする。よって$\al,\be$から連続関数
\[
\ga=\al\sqcup \be:X\sqcup Y\to [0,1]
\]
が誘導できる。$xRy\To \ga(x)=\ga(y)$であることを示そう。$R$の定義から$x\in A,y=f(x)$の時のみ調べれば良いが、
\[
\ga(x)=\al(x)=\be\circ f(x)=\be(f(x))
\]
なのですぐわかる。よって$\ga$は$R$による同値類の上で一定の値を取り、$\main=\qt{Z}{R}$なので連続関数
\[
\delta:\main\to [0,1]
\]
が誘導される。この$\delta$が$g$の拡張であることは作り方からすぐに分かる。以上で$g$は$\main$上まで連続的に拡張できるので$\main$は正規である。
\end{proof}

\begin{cor}\label{t4}
$X,Y$が$T_4$であるとき、$X\underset{f}{\cup} Y$も$T_4$である。
\end{cor}

次の命題に移る前に商空間に関する簡単な事実を述べておこう。
\begin{fact}
$E$を位相空間$S$を$E$の同値関係とし$U$を$\qt{E}{S}$の開集合とする。このとき$O=\pi^{-1}(U)$と置き、$T=S\cap (O\times O)$と置けば$T$は$O$の同値関係になり、$U$は商空間$\qt{O}{T}$に位相同型である。
\end{fact}
証明はしないが商空間の普遍性と空間$E$の開集合$P$の相対位相の開集合は$E$の開集合であることなどを使えば苦なく証明できる。
\begin{prop}
$X,Y$が遺伝的正規であるとき、$X\underset{f}{\cup} Y$も遺伝的正規である。
\end{prop}
\begin{proof}
$\main$の開集合が正規であることを示せば良い。$U$を$\main$の開集合とすると先の事実から$O=\pi^{-1}(U),S=R\cap (O\times O)$とすれば$U$は$\qt{O}{S}$に位相同型である。$\qt{O}{S}$を調べよう。\par
$P=O\cap X,Q=O\cap Y$と置くと$X,Y$の遺伝的正規性から$P,Q$は正規空間であり、$O=P\sqcup Q$である。。更に$B=O\cap A=P\cap A$とおくと之は$P$の閉集合であり、$a\in A$について$aRf(a)$なので$f(B)\subset Q$である。$g=f|B$と置こう。すると
\[
\qt{P\sqcup Q}{S}=P\underset{g}{\cup}Q\approx U
\]
であり、$P,Q$は正規であるから命題\ref{seiki}から$U$は正規である。よって$\main$は正規である。
\end{proof}

\begin{cor}
$X,Y$が$T_5$であるとき、$X\underset{f}{\cup} Y$も$T_5$である。
\end{cor}

次の命題に移る前に完全正規性に関する事実を述べよう
\begin{fact}
$X$が完全正規である事と$X$の任意の閉集合$F$上の関数$g:F\to [0,1]$をそのゼロ集合$g^{-1}(0)$を保って$X$上まで拡張できる、つまり関数$h:X\to[0,1]$が存在して
\[
h|F=g,\ h^{-1}(0)=g^{-1}(0)
\]
となるものが存在することは同値である。
\end{fact}
電波通信のどこかにこの事実の証明は載っていると思う。
\begin{prop}
$X,Y$が完全正規であるとき、$X\underset{f}{\cup} Y$も完全正規である。
\end{prop}
\begin{proof}
$\main$の任意の閉集合がゼロ集合であることを示せば良い。ほとんど命題\ref{seiki}の証明と一緒である。\par
$F$を$\main$の閉集合とする。これがゼロ集合となることを示そう。さて、$\pi^{-1}(F)$は$X\sqcup Y$の閉集合であり、
\[
K=\pi^{-1}(F)\cap X,\ L=\pi^{-1}(F)\cap Y
\]
と置けば$K,L$はそれぞれ$X,Y$の閉集合である。$Y$が完全正規であることから連続関数$\be:Y\to[0,1]$が存在して
\[
\be^{-1}(0)=L
\]
となる。
そして$X$上の閉集合$A\cup K$について$h:A\cup K\to[0,1]$を
\[
h(x)=
\begin{cases}
\be \circ f(x) & \text{$(x\in A)$}\\ 
0 & \text{$(x\in K)$}
\end{cases}
\]
と置く。之がちゃんと定義できていることを見よう。$x\in A\cap K$を取ると$xRf(x)$より$f(x)\in \pi^{-1}(F)$で$f(x)\in Y$であるから$f(x)\in L$であるから
\[
\be \circ f(x)=\be(f(x))=0
\]
なので$h$はちゃんと$A\cup K$上で定義されており、$A,K$それぞれの上で連続なので$h$は$A\cup K$上で連続である。更に定義から明らかに
\[
h^{-1}(0)=K
\]
である。
そこで$X$の完全正規性から$h$のゼロ集合を保ってを$X$上までに連続的に拡張できる。その拡張を$\al$とする。さて$\al,\be$から連続関数
\[
\ga=\al\sqcup \be:X\sqcup Y\to [0,1]
\]
が誘導できる。$xRy\To \ga(x)=\ga(y)$であることを示そう。$R$の定義から$x\in A,y=f(x)$の時のみ調べれば良いが、定義から
\[
\ga(x)=\al(x)=\be\circ f(x)=\be(f(x))
\]
なのですぐわかる。よって$\ga$は$R$による同値類の上で一定の値を取り、$\main=\qt{Z}{R}$なので連続関数
\[
\delta:\main\to [0,1]
\]
が誘導される。この$\delta$について
\[
\delta^{-1}(0)=F
\]
となることは
\[
\al^{-1}(0)=K,\ \be^{-1}(0)=L
\]
からよくわかる。よって$F$はゼロ集合なので$\main$は完全正規である。

\end{proof}

\begin{cor}
$X,Y$が$T_6$であるとき、$X\underset{f}{\cup} Y$も$T_6$である。
\end{cor}
次の命題に移る前に次の重大な２つの事実を述べよう
\begin{lem}
$K$を局所コンパクトハウスドルフ空間とし、$f:X\to Y$を等化写像とする。このとき
\[
f\times 1_K:X\times K\to Y\times K
\]
も等化写像である。
\end{lem}
\begin{proof}
$g=f\times 1_K$として
\[
\mbox{$g^{-1}(O)$が$X\times K$の開集合}\To \mbox{$O$が$Y\times K$の開集合}
\]
を示せば良い。$(a,b)\in O$を任意に与え、これが開集合であることを示す。$c\in f^{-1}(a)$を適当に選ぶと$(c,b)\in g^{-1}(O)$であり、仮定から$g^{-1}(O)$は$X\times K$の開集合なので$X$の$c$の開近傍$N$と$b$のコンパクト近傍$L$が存在して$(c,b)\in N\times L\subset g^{-1}(O)$となる。今此処で
\[
U=\{x\in X\ |\ \{x\}\times L\subset g^{-1}(O)\}
\]
と$U\subset X$を定義しよう。$L$のコンパクト性とtube lemmaから$U$は$U$は$X$の開集合となり$c\in U$となる。さて$f^{-1}(f(U))=U$を示そう。まず$U\subset f^{-1}(f(U))$はそもそも一般に成り立つ。逆を示そう。$x\in f^{-1}(f(U))$とする。すると
\[
\{x\}\times L\subset f^{-1}(f(U))\times L=f^{-1}(f(U))\times 1^{-1}_K(1_kL)=g^{-1}(g(U\times L))\subset g^{-1}(O)
\]
なので確かに$x\in U$で$f^{-1}(f(U))=U$が成立する。よって$U$が開集合であることと$f$が等化写像であることから$f(U)$は$Y$の開集合である。さらに定義から$(a,b)\in f(U)\times L\subset O$であるから$(a,b)$は$O$の内点となる。$(a,b)$は任意だったので$O$は$Y$の開集合となる。
\end{proof}

\begin{thm}
$T_{3.5}$空間$X$について以下は同値
\begin{align}
&\mbox{$X$はパラコンパクト}\\
&\mbox{任意のコンパクトハウスドルフ空間Zについて$X\times Z$は$T_4$}
\end{align}
\end{thm}
\setcounter{equation}{0}
\begin{prop}
$X,Y$がパラコンパクトハウスドルフであるとき$\main$もパラコンパクトハウスドルフである。
\end{prop}
\begin{proof}
パラコンパクトハウスドルフ空間は$T_4$なので系\ref{t4}から$\main$も$T_4$である。さて$K$を任意のコンパクトハウスドルフ空間とする。$g=f\times 1_K:A\times K\to Y\times K$として
$(X\times K)\underset{g}{\cup} (Y\times K)$を考える。$W=X\times K\sqcup Y\times K$として
\[
\{\Bigl((a,k),(f(a),k)\Bigr)\in W\times W\ |\ (a,z)\in A\times K\}
\]
から生成される最小の同値関係を$S$とすれば
\[
(X\times K)\underset{g}{\cup} (Y\times K)=\qt{W}{S}
\]
である。そして$\rho:W\to \qt{W}{S}$を商写像とする。ここで商写像$\pi:Z\to \qt{Z}{R}$について
\[
\pi\times 1_K:Z\times K\to (\qt{Z}{R})\times K
\]
は等化写像となる。つまり同値関係$T$を$xTy\overset{def}{\iff}(\pi\times 1_K)(x)=(\pi\times 1_K)(y)$と定義すれば命題\ref{syo}から
\[
\qt{(Z\times K)}{T}\approx (\qt{Z}{R})\times K
\]
となる。
さて$W=(X\times K)\sqcup (Y\times K)=(X\sqcup Y)\times K$	であるので$S,T$は同じ集合上で定義された同値関係であり、定義から明らかに$xSy\iff xTy$なので$S=T$なので
\[
(X\times K)\underset{g}{\cup} (Y\times K)=\qt{W}{S}=\qt{(Z\times K)}{T}\approx (\qt{Z}{R})\times K
\]$X,Y$がパラコンパクトであることから$X\times K,Y\times K$は$T_4$空間で系\ref{t4}から
$(X\times K)\underset{g}{\cup} (Y\times K)$も$T_4$空間となる。よってこの空間に同型な$(\qt{Z}{R})\times K$も$T_4$空間である。$\qt{Z}{R}$は$T_4$で$K$は任意のコンパクトハウスドルフ空間なので$\qt{Z}{R}$はパラコンパクトになる。よって$\qt{Z}{R}=\main$なので$\main$はパラコンパクトハウスドルフ空間である。
\end{proof}

\section{帰納極限}

\begin{df}
可算個の位相空間の族$\{X_i\}_{i\in \nn}$について各$X_i$は$X_{i+1}$の部分空間となっており、と写像$\{f_i\}_{i\in \nn},\ f_i:X_i\hk X_{i+1}$は包含写像の族であるとする。
\[
X_1\overset{f_1}{\hk} X_2  \overset{f_2}{\hk}X_3\overset{f_3}{\hk}\cdots X_n\overset{f_n}{\hk} \cdots
\]
このとき$X=\bigcup_{i\in \nn}X_i$に対して
\[
\mbox{$F$が$X$の閉集合}\overset{def}{\iff}\mbox{各$i\in \nn$について$X_i\cap F$が$X_i$の閉集合}
\]
として位相を定義する。この位相空間$X$を可算包含列$\{X_i\}_{i\in \nn}$の帰納極限位相空間と呼び、ときには$X=\varinjlim X_i$という記号で表す。\footnote{帰納極限位相空間自体はもっと一般的な状況で定義するできるが、この文書で使わないので可算包含列の帰納極限に絞ってこの様に定義した。}
また$X$の位相の定め方から
\[
\mbox{$U$が$X$の開集合}\overset{def}{\iff}\mbox{各$i\in \nn$について$X_i\cap U$が$X_i$の開集合}
\]
となる。またまた位相の定め方から
\[
\mbox{$\al;X\to Y$が連続}\overset{def}{\iff}\mbox{各$i\in \nn$について$\al:|X_i\to Y_i$が$X_i$上の連続写像}
\]
\end{df}
\begin{cau}
$X$の位相は次のように考えることも出来る。\par 
$Y_i=X\times \{i\}$として$Y_i$には$X_i$と同相となるような位相を定義し直和空間$Y=\coprod_{i\in \nn}Y_i$を作る。そして
\[
\{\Bigl((x,i),(y,j)\Bigr)\in Y\times Y\ |\ \mbox{$i<j$で$x=y$}\}
\]
から生成される最小の同値関係を$R$としたときの商空間\footnote{もちろん商位相を定義する。}$\qt{Y}{R}$は$X=\varinjlim X_i$と同相である。この事の証明は簡単だから省略する。\par
この様に見なすことによって商位相について知られていることを$\varinjlim X_i$にも使う事ができる。
\end{cau}
以下位相空間$X=\varinjlim X_i$の性質を見よう
\begin{prop}
各$X_i$が$T_1$ならば$X=\varinjlim X_i$も$T_1$である。
\end{prop}

\begin{prop}\label{seiki2}
可算包含列$\{X_i\}_{i\in \nn}$について各$i\in \nn$で$X_i$が$X_{i+1}$の閉集合になっているとする。このとき
各$X_i$が正規ならば$X=\varinjlim X_i$も正規である。
\end{prop}
\begin{proof}
$X$上の任意の閉集合$F$と$F$上の連続関数$g:F\to [0,1]$が$X$全体に拡張出来る事を示せば良い。いま帰納的に関数列$\{h_i\}_{i\in \nn}$を$h_i:X_i\to [0,1]$で$h_{k+1}$は$h_k$の$X_{k+1}$への拡張で
\[
h_j|F\cap X_j=g|F\cap X_j
\]
を充たすようにを構成しよう。\par
まず$i=1$のとき：$X$の位相の定め方から$F\cap X_1$は$X_i$の閉集合である。そこで$X_1$の正規性から$F\cap X_1$上の関数$g|F\cap X_1:F\cap X_1\to[0,1]$を$X_1$まで連続的に拡張出来る。これを$h_1$としよう。すると確かに
\[
h_1|F\cap X_1=g|F\cap X_1
\]は成り立っている。\par
$i=k$まで構成されているとして$i=k+1$のとき：$X$の位相の定め方から$F\cap X_{k+1}$は$X_{k+1}$の閉集合になっている。いま$X_{k+1}$の閉集合$X_k\cup (F\cap X_{k+1})$上の関数$v$を
\[
v(x)=
\begin{cases}
h_k(x) & \text{$(x\in X_k)$}\\
g|F\cap X_{k+1} & \text{$(x\in F\cap X_{k+1})$}
\end{cases}
\]
と定義する。$h_k$と$g|F\cap X_{k+1}$は$X_k\cap F\cap X_{k+1}=X_k\cap F$上同じ値をとるので確かに$v$は定義できて、$X_k$上でも$F\cap X_{k+1}$でも連続なので$v:X_k\cup (F\cap X_{k+1})\to [0,1]$は連続である。そして$X_{k+1}$の正規性から$v$を$X_{k+1}$までに拡張した関数を$h_{k+1}$とすれば確かに$h_{k+1}$は$h_k$の拡張で
\[
h_{k+1}|F\cap X_{k+1}=g|F\cap X_{k+1}
\]
が成り立っている。\par
以上から関数列$\{h_i\}_{i\in \nn}$が構成できた。そこで
\[
H:X\to [0,1]
\]
を
\[
\mbox{$x\in X_i$のとき$H(x)=h_i(x)$}
\]
と定義するとこれはちゃんと定義できて、$H|X_i=h_i$で$h_i$は$X_i$上連続なので$X$の位相の定め方から$H$は連続関数であり、関数列$\{h_i\}_{i\in \nn}$の定め方から$g:F\to [0,1]$の拡張になっている。よって$X=\varinjlim X_i$は正規である。
\end{proof}

%%%%%%%%%%%%%%遺伝的正規
\begin{lem}\label{open}
$X=\varinjlim X_i$の開集合$U$について$U$は可算包含列$\{X_i\cap U\}_{i\in \nn}$の帰納極限になっている。つまり
\[
U=\varinjlim (X_i\cap U)
\]
\end{lem}
\begin{proof}$U$の部分集合$O$について
\[
\mbox{各$i\in \nn$について$(X_i\cap U)\cap O=X_i\cap O$が$X_i\cap U$の開集合}\To \mbox{$O$は$X$の開集合}
\]
を示せば良い。$X_i\cap U$は$X_i$の開集合であり、$X_i\cap O$は$X_i\cap U$の開集合である。よって$X_i\cap O$は$X_i$の開集合である。（開集合の相対開集合は元の空間の開集合。）よって
\[
\mbox{各$i\in \nn$について$X_i\cap O$は$X_i$の開集合}
\]
が成り立つので$X$の位相の定め方から$O$は$X$の開集合となり、特に$U$の開集合になる。
\end{proof}

\begin{prop}
可算包含列$\{X_i\}_{i\in \nn}$について各$i\in \nn$で$X_i$が$X_{i+1}$の閉集合になっているとする。このとき
各$X_i$が遺伝的正規ならば$X=\varinjlim X_i$も遺伝的正規である。
\end{prop}
\begin{proof}
$X$の任意の開集合$U$が正規であることを示せばよい。\par
さて$X$の位相の定義から$U\cap X_i$は$X_i$の開集合となるが$X_i$は遺伝的正規なので$U\cap X_i$は正規空間である。そして$X_i$が$X_{i+1}$の閉集合となっているから$U\cap X_i$は$U\cap X_{i+1}$の閉集合となっている。
よって補題\ref{open}と命題\ref{seiki2}から$U=\varinjlim (U\cap X_i)$は正規空間になる。
\end{proof}

%%%%%%%%%%%完全正規
\begin{prop}
可算包含列$\{X_i\}_{i\in \nn}$について各$i\in \nn$で$X_i$が$X_{i+1}$の閉集合になっているとする。このとき
各$X_i$が完全正規ならば$X=\varinjlim X_i$も完全正規である。
\end{prop}
\begin{proof}
命題\ref{seiki2}とほとんど同様である。\par
$X$上の任意の閉集合$F$が$X$のゼロ集合であることを示せば良い。いま帰納的に関数列$\{h_i\}_{i\in \nn}$を$h_i:X_i\to [0,1]$で$h_{k+1}$は$h_k$の$X_{k+1}$への拡張で
\[
h_j^{-1}(0)=X_j\cap F
\]
を充たすようにを構成しよう。\par
まず$i=1$のとき：$X$の位相の定め方から$F\cap X_1$は$X_i$の閉集合である。そこで$X_1$の完全正規性から連続関数$h_1:X_1\to [0,1]$が存在して
\[
h_1^{-1}(0)=X_i\cap F
\]
が成り立つので良い。\par
$i=k$まで構成されているとして$i=k+1$のとき：$X$の位相の定め方から$F\cap X_{k+1}$は$X_{k+1}$の閉集合になっている。いま$X_{k+1}$の閉集合$X_k\cup (F\cap X_{k+1})$上の関数$v$を
\[
v(x)=
\begin{cases}
h_k(x) & \text{$(x\in X_k)$}\\
0 & \text{$(x\in F\cap X_{k+1})$}
\end{cases}
\]
と定義する。$h_k^{-1}(0)=X_k\cap F$なので$h_k$は$X_k\cap F\cap X_{k+1}=X_k\cap F$上$0$なので確かに$v$は定義できて、$X_k$上でも$F\cap X_{k+1}$でも連続なので$v:X_k\cup (F\cap X_{k+1})\to [0,1]$は連続である。更に
\[
v^{-1}(0)=X_{k+1}\cap F
\]
である。
そして$X_{k+1}$の完全正規性から$v$を$X_{k+1}$まで$v$のゼロ集合を保ったまま拡張した関数を$h_{k+1}$とすれば確かに$h_{k+1}$は$h_k$の拡張で
\[
h_{k+1}^{-1}(0)=X_{k+1}\cap F
\]
が成り立っている。\par
以上から関数列$\{h_i\}_{i\in \nn}$が構成できた。そこで
\[
H:X\to [0,1]
\]
を
\[
\mbox{$x\in X_i$のとき$H(x)=h_i(x)$}
\]
と定義するとこれはちゃんと定義できて、$H|X_i=h_i$で$h_i$は$X_i$上連続なので$X$の位相の定め方から$H$は連続関数であり、関数列$\{h_i\}_{i\in \nn}$の定め方から
\[
H^{-1}(0)=\bigcup_{i\in\nn}(F\cap X_i)=F
\]
となるので$F$はゼロ集合である。よって$X=\varinjlim X_i$は完全正規である。

\end{proof}
%%%%%%%%%%パラコンパクト
\begin{lem}\label{morita}
$K$を局所コンパクトハウスドルフ空間すると$X\times K$は可算包含列$\{X_i\times K\}_{i\in \nn}$の帰納極限になっているつまり
\[
X\times K=\varinjlim (X_i\times K)
\]
が成り立つ。もちろん$\{X_i\times K\}_{i\in \nn}$の包含写像は$f_i\times 1_K$である。
\end{lem}
\begin{proof}$X\times K$の部分集合$U$について
\[
\mbox{各$i\in \nn$について$(X_i\times K)\cap U$が$X_i\times K$の開集合}\To\mbox{$U$が$X\times K$の開集合}
\]
を示せばよい。$(a,b)\in U$を任意に与える。そして$x\in X_n$とする。
\[
H=\{y\in K\ |\ (a,y)\in O\}
\]
と$H\subset Y$を定義すると之は開集合である。なぜなら
\[
H=\{y\in K\ |\ (a,y)\in O\cap (X_n\times K)\}
\]
であり、$O\cap (X_n\times K)$は$X_n$の開集合だからである。さて明らかに$b\in H$で$K$は局所コンパクトハウスドルフ空間なので$b$のコンパクト近傍$L$が存在して$b\in L\subset H$となる。そこで
\[
U=\{x\in X\ |\ \{x\}\times L\subset O\}
\]
と$U\subset X$と定義する。このとき各$i\in \nn$について
\[
U\cap X_i=\{x\in X\ |\ \{x\}\times L\subset O\cap(X_i\times K)\}
\]
であり、$O\cap(X_i\times K)$が$X_i\times K$の開集合で$L$がコンパクトなのでtube lemmaから$U\cap X_i$は$X_i$の開集合である。よって$X$の位相の定め方から$U$は$X$の開集合である。そして$a\in U$である。よって$U$の定義から
\[
(a,b)\in U\times L\subset O
\]
で$U\times L$は$(a,b)$の近傍になるから$(a,b)$は$O$の内点である。$(a,b)$は任意だったから$O$は開集合である。
\end{proof}
\begin{prop}
可算包含列$\{X_i\}_{i\in \nn}$について各$i\in \nn$で$X_i$が$X_{i+1}$の閉集合になっているとする。このとき
各$X_i$がパラコンパクトハウスドルフならば$X=\varinjlim X_i$もパラコンパクトハウスドルフである。
\end{prop}
\begin{proof}
まずパラコンパクトハウスドルフ空間は$T_4$なので$X$は$T_4$となる。いま任意にコンパクトハウスドルフ空間$K$を与えると$X_i$がパラコンパクトハウスドルフであることから$X_i\times K$は$T_4$である。よって命題\ref{seiki2}と補題\ref{morita}から
\[
X\times K=\varinjlim (X_i\times K)
\]
も$T_4$である。よって$X$は正規で任意のコンパクトハウスドルフ空間$K$について$X\times K$が$T_4$なので$X=\varinjlim X_i$もパラコンパクトである。
\end{proof}

\section{CW複体}
このセクシオンに於いて$D^m$で$m$次元閉球を表し$\pt D^m$でその境界を表す。つまり
\[
D^m=\{x\in\rr^m\ |\ \|x\|\le 1\},\ \pt D^m=\{x\in \rr^m\ |\ \|x\|=1\}
\]
である。また、$D^0$は１点であり、$\pt D^0=\O$とする。また$\zz_{\ge-1}=\{n\in \zz\ |\ n\ge-1\}$とする。
\begin{df}[相対CW複体とCW複体]\footnote{このCW複体の定義とホワイトヘッドが定義した所謂条件$(C),(W)$を充たすハウスドルフ空間というCW複体の定義が一致することは例えば\cite{Hat}のAppendixを参照されたし}
以下を充たす可算包含列$\{X^{i}\}_{i\in\zz_{\ge-1}}$の帰納極限位相空間であるような空間を位相空間$A$を核とした相対CW複体$X=\varinjlim X^{(i)}$という。また、$X$と$A$を組にして$(X,A)$を相対CW複体と言ったりする。$A=\O$のとき$(X、\O)$を単にCW複体と言う。\par
$(I){}$$
X^{(-1)}=A$\par
$(I\!I){}$各$k\in \zz_{\ge-1}$について
写像の族$\{\varphi_{\al}^{k+1}:\pt D^{k+1}\to X^{(k)}\}_{\al\in I_{k+1}}$が存在して
\[
\Phi^{k+1}=\coprod_{\al\in I_{k+1}}\varphi_{\al}^{k+1}:\coprod_{\al\in I_{k+1}}\pt D^{k+1}\to X^{(k)}
\]
と置いて、$\dpst \coprod_{\al\in I_{k+1}}\pt D^{k+1}\subset \coprod_{\al\in I_{k+1}}D^{k+1}$として
\[
X^{(k+1)}=\left(\coprod_{\al\in I_{k+1}}D^{k+1}\right)\underset{\Phi^{k+1}}{\cup}X^{(k)}
\]
となっている。つまり$X^{(k+1)}$は$\dpst \coprod_{\al\in I_{k+1}}D^{k+1}$を$\Phi^{k+1}$で$X^{(k)}$に接着した空間である。
\[\xymatrix{
& \coprod_{\al\in I_{k+1}}\pt D^{k+1} \ar@{^{(}->}[r]^{\iota}\ar@{->}[d]_{\Phi^{k+1}}& \coprod_{\al\in I_{k+1}}D^{k+1} \ar@{->}[d]&\\ 
& X^{(k)} \ar@{->}[r]_{j_{k}}& X^{(k+1)} &
}\]
\end{df}
上の定義において$j_k:X^{(k)}\to X^{(k+1)}$は命題\ref{Y}より閉写像で埋め込みになっているので$X^{(k)}$は$X^{(k+1)}$の閉部分空間とみなせる。このこととCW複体の定義とセクション1,2次の系は明らかである。
\begin{cor}
$A$が$T_1$ならば$A$を核とした相対CW複体$(X,A)$も$T_1$
\end{cor}


\begin{cor}
$A$が正規ならば$A$を核とした相対CW複体$(X,A)$も正規
\end{cor}

\begin{cor}
$A$が遺伝的正規ならば$A$を核とした相対CW複体$(X,A)$も遺伝的正規
\end{cor}

\begin{cor}
$A$が完全正規ならば$A$を核とした相対CW複体$(X,A)$も完全正規
\end{cor}

\begin{cor}
$A$がパラコンパクトハウスドルフならば$A$を核とした相対CW複体$(X,A)$もパラコンパクトハウスドルフ
\end{cor}

\begin{cor}
CW複体はパラコンパクトハウスドルフ$T_6$空間である。
\end{cor}
	
\begin{thebibliography}{9}
\bibitem{1}J.Dugundji,{\it Topology},William C Brown Pub ,1966
\bibitem{Hat}A.Hatcer,{\it Algebraic topology},Cambridge univ. press, 2002 
\bibitem{4}K.Morita,{\it On spaces having the weak topology with respect to closed coverings},Proc.Japan Acad.Vol29 No.10(1953)
\bibitem{pear}A.R.Pears,{\it DIMENSION THEORY OF GENERAL SPACES},Cambridge Univ.Press,1975
\end{thebibliography}




			\end{document}



\begin{comment}
[可換図式]
\[\xymatrix{
&   & (2) \ar@{=>}[rd]&  &  &  & (5) \ar@{=>}[rd]&\\ 
& (1)\ar@{=>}[ru] \ar@{=>}[rd]&  &(4)\ar@{=>}[r]&(7)& (1)\ar@{=>}[ru] \ar@{=>}[rd]& & (7)&\\ 
&   & (3) \ar@{=>}[ur]&  &  &  &(6) \ar@{=>}[ur]&
}\]
[\bigin \endを使わない方法]

{\prop[aa] aaa}
で
\begin{prop}[aa]
aaa
\end{prop}
と同じ\df \thm \proof でも同様

[直和]
$\amalg$

[同値関係でよく使う波線]
\sim

[引数付きの新命令]
\newcommand{\prn}[1]{\|#1\|_p}

[番号のリセット]

\setcounter{equation}{0}


[字体の変更]

{\rm hogehoge}と使う。

\rm ローマン
\it イタリック
\sl 斜体

[supp]

\newcommand{\supp}{\text{{\rm supp}}}

[中央揃えの改行]

	\begin{eqnarray*}
		hoge1&=&hogehoge
		hoge2&=&hogehofe

	\end{eqnarray*}

&&の部分で揃えられる。






[場合分け]

	\begin{cases}
		hoge1 & \text{状況１}\\
		hoge2 & \text{状況２}\\
		・
		・
		・
	\end{cases}



\end{comment}





